\documentclass{article} 
\usepackage[a4paper, margin=3cm]{geometry}

\usepackage{cite}

\usepackage{pgfgantt}
\usepackage{pdflscape}


\title{\textbf{Network Influence Analysis Framework}} 
\author{Antoni Zyla}
\date{21/07/2025}

\begin{document}

\maketitle

\section{Abstract}

To create a framework for the analysis of influence within networks,
particularly the investigation of heuristics for finding the most
influential nodes in an efficient manner.

\section{Problem Discussion}

Network Influence Analysis relates to studying social networks
to determine the set of $k$ nodes which have the most 'influence',
or greatest spread of information through the network, in the
network. Using brute force is incredibly time consuming
on the order of $O \left( |V| \choose k \right)$. Using certain heuristics
such as the degree of a node to choose the seed nodes is significantly
more efficient running in polynomial time, I will focus on allowing for
simple implementations and analysis of different heuristics.

\vspace*{12pt}

There is a significant amount of information out there regarding how
these models are implemented and computed, there are graph packages in most
programmming languages which make it relatively simple to implement them
as described.

\vspace*{12pt}

A typical piece of research in this area would be to attempt to find a
new heuristic to find the set of seed nodes which have a greater influence
than other heuristics. To do this the researcher would need to implement
the model, implement representing their chosen networks in their model,
implement the known heuristics and their new hypothesised heuristic and
analyse the influence spread of each of them to be able to see if there is a
measurable difference.

\vspace*{12pt}

Creation of a framework for performing these tasks would allow for creating
such a study to be simpler, quicker and more reproducible due to a consistent
environment.

\section{Objectives}

\begin{itemize}
	\item Import and Export of different graph formats
	\item Generation of synthetic graphs (\textbf{Kleinberg, Eppstein})
	\item Independent Cascade Model (\textbf{IC})
	\item Linear Threshold Model (\textbf{LTM})
	\item Decreasing Cascade Model (\textbf{DC})
	\item Different Sampling of large graphs to make them manageable
	\item Different heuristics to find seed nodes
	\item Simple Creation of new sampling methods
	\item Simple Creation of new heuristics
	\item Python \textbf{whl} and \textbf{tarball} packaging to allow easy installation
\end{itemize}

\section{Methodology}

The project will be broken up into three main stages; research, implementation and
evaluation.

\vspace*{12pt}

The research stage will require a detailed literature review of articles
pertaining to network influence modelling, methods, heuristics and sampling.
I will identify the validity and applications of each source and find if there
are any gaps in the presented knowledge to research further.

\vspace*{12pt}

The implementation stage of the project will require a methodology that can
be adapted to changing requirements or a shift in focus based on information
learnt later in the research phase. I will use an agile-like methodology, of
one to two week segments determined by meetings with my supervisor with any
changes to requirements determined during these meetings.

\vspace*{12pt}

The final stage will be to use my framework to replicate the results of the
researched papers, showcase how to create new sampling methods and heuristics
and how to compare the results of different sampling, heuristics and models
on different networks.


\begin{landscape}

	\section{Timelines}

  \hspace{1cm}

  \newganttlinktype{end-start}{
  \ganttsetendanchor{on left=0}
  \ganttsetstartanchor{on right=0.5}
  \draw[/pgfgantt/link,->] (\xLeft,\yUpper) -| (\xRight,\yLower);
  }

  \newganttlinktype{start-start}{
  \ganttsetendanchor{on left=0}
  \ganttsetstartanchor{on left=0.5}
  \draw[/pgfgantt/link,->] (\xLeft,\yUpper) -| (\xRight,\yLower);
  }

  \begin{ganttchart}[
    x unit=0.7cm,
    y unit title=0.5cm,
    y unit chart=0.5cm,
    hgrid=true,
    vgrid=true,
    title label anchor/.style={below=-1.5ex},
    title height=1,
    incomplete/.style={fill=white},
    progress label text={},
    bar height=0.6,
]{1}{32}
  \gantttitle{\textbf{Term 1}}{10}
  \gantttitle{\textbf{Christmas}}{4}
  \gantttitle{\textbf{Term 2}}{10}
  \gantttitle{\textbf{Easter}}{5}
  \gantttitle{\textbf{Term 3}}{3}\\
  \gantttitle{W1}{1}
  \gantttitle{W2}{1}
  \gantttitle{W3}{1}
  \gantttitle{W4}{1}
  \gantttitle{W5}{1}
  \gantttitle{W6}{1}
  \gantttitle{W7}{1}
  \gantttitle{W8}{1}
  \gantttitle{W9}{1}
  \gantttitle{W10}{1}%term 1 done
  \gantttitle{W1}{1}
  \gantttitle{W2}{1}
  \gantttitle{W3}{1}
  \gantttitle{W4}{1}%christmas done
  \gantttitle{W1}{1}
  \gantttitle{W2}{1}
  \gantttitle{W3}{1}
  \gantttitle{W4}{1}
  \gantttitle{W5}{1}
  \gantttitle{W6}{1}
  \gantttitle{W7}{1}
  \gantttitle{W8}{1}
  \gantttitle{W9}{1}
  \gantttitle{W10}{1}%term 2 done
  \gantttitle{W1}{1}
  \gantttitle{W2}{1}
  \gantttitle{W3}{1}
  \gantttitle{W4}{1}
  \gantttitle{W5}{1}% Easter Done
  \gantttitle{W1}{1}
  \gantttitle{W2}{1}
  \gantttitle{W3}{1} \\ %t3 done

  \ganttbar[name=spec]{\textbf{Specification}}{1}{1}\\    
  \ganttbar[name=progress]{\textbf{Progress Report}}{4}{5} \\ 
  \ganttbar[name=literature]{Literature Review}{2}{2}\\
  \ganttbar[name=architecture]{Architecture \& Definition}{3}{4} \\
  \ganttbar[name=implement]{Implementation}{6}{14}\\
  \ganttbar[name=present]{\textbf{Presentation}}{14}{17} \\
  \ganttbar[name=report]{\textbf{Final Report}}{18}{30}

  % arrows

  \ganttlink[]{spec}{progress}
  \ganttlink[]{spec}{literature}
  \ganttlink[]{literature}{architecture}
  \ganttlink[]{progress}{implement}

\end{ganttchart}

\end{landscape}

\section{Resources and Risks}

\section{Legal, Social, Ethical and Professional Considerations}


\bibliography{sources}{}
\bibliographystyle{plain}

\end{document}
